% GENERATED FILE. Edit canonical lessons, then run npm run generate:books.
% canonical-source-hash: fnv1a64:661d9f2770f2d3d0
% canonical-lessons: GE-C436-gasse, GE-C436-sackgasse, GE-C436-allee, GE-C436-einbahnstrasse, GE-C436-passant

\chapter{Kleine und große Straßen}
\label{ch:ge-kleine-und-gro-e-stra-en}
\hlchaptermodality{436}

\begin{chapteropening}
\textbf{By the end of this chapter:} \emph{I can describe streets: the lane, the dead end, the avenue, the one-way street and the passer-by.}

Complete the last lesson of chapter 436: I can describe streets: the lane, the dead end, the avenue, the one-way street and the passer-by.
\end{chapteropening}

\section[die Gasse]{die Gasse — the lane, the alley}

\label{lesson:GE-C436-gasse}



\begin{quote}
Before the new one: say the German for to turn back, then the German for to get lost driving.
\end{quote}

\subsection*{You'll want to know: die Gasse}
\textbf{die Gasse} — \textquotedblleft{}the lane, the alley\textquotedblright{}.

\begin{quote}
Die Gassen im alten Teil der Stadt sind eng.
\end{quote}

\begin{quote}
\emph{The lanes in the old part of town are narrow.}
\end{quote}

Plural: \emph{die Gassen}. In Austria, \emph{die Gasse} is an ordinary word for a street.

\begin{cousinweb}
\emph{Gasse} is inherited from Old High German \emph{gazza} and Proto-Germanic \emph{gatwōn} (EtymDB). Old Norse \emph{gata} 'road' is usually counted as a cousin.
\end{cousinweb}

\begin{grammarlens}[title={What you've built}]
One more word for telling, explaining and describing at B1.
\end{grammarlens}

\subsection*{Your turn}
\begin{itemize}
\raggedright
  \item \emph{Say it:} \emph{die Gasse}
  \item \emph{Say it:} \emph{die Gasse}, once more
  \item \emph{Say it:} a sentence of your own with \emph{die Gasse}
  \item \emph{Recall:} say the German for to turn back, then the German for to get lost driving, then say \emph{die Gasse} again
\end{itemize}

\begin{tcolorbox}[breakable,colback=teal!4,colframe=teal!35!black,title={Before you move on}]
What does die Gasse mean? (\textquotedblleft{}The lane, the alley\textquotedblright{}.) Say it once more.
\end{tcolorbox}
\section[die Sackgasse]{die Sackgasse — the dead end, the cul-de-sac}

\label{lesson:GE-C436-sackgasse}



\begin{quote}
Before the new one: say the German for to get lost driving, then the German for the lane.
\end{quote}

\subsection*{You'll want to know: die Sackgasse}
\textbf{die Sackgasse} — \textquotedblleft{}the dead end, the cul-de-sac\textquotedblright{}.

\begin{quote}
Die Straße ist eine Sackgasse, wir müssen zurückfahren.
\end{quote}

\begin{quote}
\emph{The street is a dead end; we have to drive back.}
\end{quote}

Plural: \emph{die Sackgassen}. Also figurative: a situation with no way out.

\begin{cousinweb}
\emph{Sackgasse} is built from \emph{Sack} 'sack', which goes back to Latin \emph{saccus} (EtymDB; Lewis \& Short), and \emph{Gasse} 'lane': a lane closed at one end like a bag.
\end{cousinweb}

\begin{grammarlens}[title={What you've built}]
One more word for telling, explaining and describing at B1.
\end{grammarlens}

\subsection*{Your turn}
\begin{itemize}
\raggedright
  \item \emph{Say it:} \emph{die Sackgasse}
  \item \emph{Say it:} \emph{die Sackgasse}, once more
  \item \emph{Say it:} a sentence of your own with \emph{die Sackgasse}
  \item \emph{Recall:} say the German for to get lost driving, then the German for the lane, then say \emph{die Sackgasse} again
\end{itemize}

\begin{tcolorbox}[breakable,colback=teal!4,colframe=teal!35!black,title={Before you move on}]
What does die Sackgasse mean? (\textquotedblleft{}The dead end, the cul-de-sac\textquotedblright{}.) Say it once more.
\end{tcolorbox}
\section[die Allee]{die Allee — the avenue (lined with trees)}

\label{lesson:GE-C436-allee}



\begin{quote}
Before the new one: say the German for the lane, then the German for the dead end.
\end{quote}

\subsection*{You'll want to know: die Allee}
\textbf{die Allee} — \textquotedblleft{}the avenue (lined with trees)\textquotedblright{}.

\begin{quote}
Vom Schloss führt eine lange Allee zum Dorf.
\end{quote}

\begin{quote}
\emph{A long avenue leads from the palace to the village.}
\end{quote}

Plural: \emph{die Alleen}.

\begin{cousinweb}
\emph{Allee} was borrowed from French \emph{allée} 'path, walk' (EtymDB), from \emph{aller} 'to go'.
\end{cousinweb}

\begin{grammarlens}[title={What you've built}]
One more word for telling, explaining and describing at B1.
\end{grammarlens}

\subsection*{Your turn}
\begin{itemize}
\raggedright
  \item \emph{Say it:} \emph{die Allee}
  \item \emph{Say it:} \emph{die Allee}, once more
  \item \emph{Say it:} a sentence of your own with \emph{die Allee}
  \item \emph{Recall:} say the German for the lane, then the German for the dead end, then say \emph{die Allee} again
\end{itemize}

\begin{tcolorbox}[breakable,colback=teal!4,colframe=teal!35!black,title={Before you move on}]
What does die Allee mean? (\textquotedblleft{}The avenue (lined with trees)\textquotedblright{}.) Say it once more.
\end{tcolorbox}
\section[die Einbahnstraße]{die Einbahnstraße — the one-way street}

\label{lesson:GE-C436-einbahnstrasse}



\begin{quote}
Before the new one: say the German for the dead end, then the German for the avenue.
\end{quote}

\subsection*{You'll want to know: die Einbahnstraße}
\textbf{die Einbahnstraße} — \textquotedblleft{}the one-way street\textquotedblright{}.

\begin{quote}
In der Einbahnstraße fahren alle Autos in eine Richtung.
\end{quote}

\begin{quote}
\emph{In the one-way street all the cars go in one direction.}
\end{quote}

Plural: \emph{die Einbahnstraßen}.

\begin{cousinweb}
\emph{Einbahnstraße} is built from \emph{ein} 'one', \emph{Bahn} 'track, lane' and \emph{Straße} 'street', which goes back to Latin \emph{(via) strāta} 'paved road' (Lewis \& Short).
\end{cousinweb}

\begin{grammarlens}[title={What you've built}]
One more word for telling, explaining and describing at B1.
\end{grammarlens}

\subsection*{Your turn}
\begin{itemize}
\raggedright
  \item \emph{Say it:} \emph{die Einbahnstraße}
  \item \emph{Say it:} \emph{die Einbahnstraße}, once more
  \item \emph{Say it:} a sentence of your own with \emph{die Einbahnstraße}
  \item \emph{Recall:} say the German for the dead end, then the German for the avenue, then say \emph{die Einbahnstraße} again
\end{itemize}

\begin{tcolorbox}[breakable,colback=teal!4,colframe=teal!35!black,title={Before you move on}]
What does die Einbahnstraße mean? (\textquotedblleft{}The one-way street\textquotedblright{}.) Say it once more.
\end{tcolorbox}
\section[der Passant]{der Passant — the passer-by}

\label{lesson:GE-C436-passant}



\begin{quote}
Before the new one: say the German for the avenue, then the German for the one-way street.
\end{quote}

\subsection*{You'll want to know: der Passant}
\textbf{der Passant} — \textquotedblleft{}the passer-by\textquotedblright{}.

\begin{quote}
Ein Passant zeigt uns den Weg zum Museum.
\end{quote}

\begin{quote}
\emph{A passer-by shows us the way to the museum.}
\end{quote}

Weak noun: \emph{den Passanten}. A woman: \emph{die Passantin}.

\begin{cousinweb}
\emph{Passant} was borrowed from French \emph{passant} 'passing (by)' (EtymDB), from \emph{passer} 'to pass', which goes back to Latin \emph{passus} 'step'.
\end{cousinweb}

\begin{grammarlens}[title={What you've built}]
One more word for telling, explaining and describing at B1.
\end{grammarlens}

\subsection*{Your turn}
\begin{itemize}
\raggedright
  \item \emph{Say it:} \emph{der Passant}
  \item \emph{Say it:} \emph{der Passant}, once more
  \item \emph{Say it:} a sentence of your own with \emph{der Passant}
  \item \emph{Recall:} say the German for the avenue, then the German for the one-way street, then say \emph{der Passant} again
\end{itemize}

\begin{tcolorbox}[breakable,colback=teal!4,colframe=teal!35!black,title={Before you move on}]
What does der Passant mean? (\textquotedblleft{}The passer-by\textquotedblright{}.) Say it once more.
\end{tcolorbox}
